\documentclass{article}
\usepackage[T1]{fontenc}
\usepackage{lmodern}
\usepackage{graphicx}
\usepackage{amsmath,amssymb}
\usepackage[a4paper,margin=2cm]{geometry}
\usepackage[absolute,overlay]{textpos}
\setlength{\TPHorizModule}{1mm}
\setlength{\TPVertModule}{1mm}
\usepackage[indonesian]{babel}
\usepackage{hyperref}
\setlength{\emergencystretch}{2em}
% unit: Halaman 66

\begin{document}

% === Halaman 66 (llm) ===

\begin{itemize}
    \item Setiap register didetak dalam fase $stop/go$ dengan menggunakan kaidah mayoritas:
    
    \textcolor{red}{"Pada setiap putaran ($i=1..64$), bit-bit tengah setiap register diperiksa dan bit mayoritasnya ($50\% + 1$) ditentukan. Jika bit tengah sebuah register sama dengan bit mayoritas, maka register tersebut didetak"}
    
    \item Biasanya pada setiap putaran dua atau tiga buah register didetak. Peluang sebuah register didetak pada setiap putaran adalah $\frac{3}{4}$.
    
    \item Ketika sebuah register didetak, semua bit detaknya di-XOR-kan dan hasilnya diletakkan pada posisi LSB (posisi ke-0) dengan mekanisme pergeseran bit-bit ke kiri.
    
    \item Bit paling kiri (MSB) terlempar keluar. Bit yang terlempar dari setiap register di-XOR-kan bersama, bit inilah yang menjadi luaran dari ketiga buah register tadi.
\end{itemize}

\end{document}
